\documentclass[12pt, a4paper]{article}

% ---------- پکیج‌های پایه (قبل از xepersian) ----------
\usepackage{xcolor}
\usepackage{geometry}
\usepackage{hyperref}
\usepackage{titlesec}
\usepackage{enumitem}
\usepackage{longtable}
\usepackage{booktabs}
\usepackage{array}
\usepackage{ragged2e}
\usepackage{fancyhdr}

% ---------- xepersian (باید آخرین پکیج باشد) ----------
\usepackage{xepersian}
\settextfont{Vazirmatn}
\setlatintextfont{Times New Roman}
\setdigitfont{Vazirmatn}

\geometry{margin=2.5cm}

\hypersetup{
    colorlinks=true,
    linkcolor=blue,
    urlcolor=blue,
    citecolor=blue
}

\titleformat{\section}{\Large\bfseries\color{blue!60!black}}{\thesection}{1em}{}
\titleformat{\subsection}{\large\bfseries}{\thesubsection}{1em}{}
\titleformat{\subsubsection}{\normalsize\bfseries}{\thesubsubsection}{1em}{}

\setlist{nosep, leftmargin=*}

\pagestyle{fancy}
\fancyhf{}
\fancyhead[R]{\small سند پرسوناها — هوم‌استور}
\fancyhead[L]{\small نسخهٔ ۱.۰}
\fancyfoot[C]{\thepage}
\renewcommand{\headrulewidth}{0.4pt}

\linespread{1.4}

\title{
    \vspace{-1cm}
    \textbf{\Huge سند پرسوناها}\\[0.3cm]
    \textbf{\LARGE پروژهٔ هوم‌استور}\\[0.5cm]
    \large کاربران هدف و ویژگی‌های آن‌ها
}
\author{
    تهیه‌کننده: ابراهیم فضلی\\
    دانشگاه: أزاد واحد زنجان\\
    ترم: ۱۴۰۴
}
\date{\today}

% ============================================================
\begin{document}

\maketitle
\thispagestyle{empty}

\vspace{1cm}

\begin{center}
\begin{tabular}{|l|l|}
\hline
\textbf{نام پروژه} & هوم‌استور (HomeStore) \\
\hline
\textbf{نسخهٔ سند} & ۱.۰ \\
\hline
\textbf{تاریخ} & \today \\
\hline
\textbf{تهیه‌کننده} & ابراهیم فضلی \\
\hline
\textbf{وضعیت} & پیش‌نویس \\
\hline
\end{tabular}
\end{center}

\newpage
\tableofcontents
\newpage

% ============================================================
\section{مقدمه}

\subsection{هدف سند}
این سند پرسوناهای کاربران هدف پروژهٔ \textbf{هوم‌استور} را تعریف می‌کند.
هدف آن ایجاد درک عمیق از نیازها، اهداف و دردهای کاربران است تا
تصمیم‌گیری‌های طراحی و توسعه بر پایهٔ آن‌ها انجام شود.

\subsection{چرا پرسونا؟}
پرسونا یک \textbf{نمایندهٔ فرضی} از یک گروه کاربری است که با داده‌های واقعی
ساخته می‌شود. پرسونا کمک می‌کند:
\begin{itemize}
    \item تیم توسعه بداند برای \textbf{چه کسی} می‌سازد.
    \item از «طراحی برای همه» پرهیز شود.
    \item اولویت‌بندی ویژگی‌ها بر پایهٔ نیاز واقعی کاربران باشد.
\end{itemize}

\subsection{روش‌شناسی}
پرسوناها بر پایهٔ موارد زیر ساخته شده‌اند:
\begin{itemize}
    \item مصاحبهٔ فرضی با مشتری
    \item مشاهدهٔ رفتار خرید آنلاین لوازم خانگی
    \item مطالعهٔ سایت‌های مشابه
    \item تجربهٔ شخصی
\end{itemize}

% ============================================================
\section{پرسونا ۱ — خانم رضایی (ادمین)}

\subsection{مشخصات دموگرافیک}
\begin{center}
\begin{tabular}{|l|l|}
\hline
\textbf{ویژگی} & \textbf{مقدار} \\
\hline
نام & فاطمه رضایی \\
\hline
سن & ۴۲ سال \\
\hline
جنسیت & زن \\
\hline
شغل & صاحب فروشگاه لوازم خانگی \\
\hline
محل زندگی & کاشان \\
\hline
تحصیلات & دیپلم \\
\hline
مهارت فنی & متوسط \\
\hline
وضعیت تأهل & متأهل، دو فرزند \\
\hline
\end{tabular}
\end{center}

\subsection{نقل‌قول}
\begin{quote}
«محصولاتم اصل و با گارانتی‌اند، اما فقط مشتری‌های محلی من را می‌شناسند.»
\end{quote}

\subsection{اهداف}
\begin{itemize}
    \item افزایش فروش فراتر از مشتریان محلی
    \item مدیریت آسان محصولات و سفارش‌ها
    \item دسترسی به مشتریان شهرهای اطراف
    \item کاهش وابستگی به فروش حضوری
\end{itemize}

\subsection{دردها}
\begin{itemize}
    \item نمی‌تواند محصولات را آنلاین نمایش دهد
    \item سفارش‌ها تلفنی و پراکنده‌اند
    \item موجودی را دستی ثبت می‌کند
    \item رقابت با فروشگاه‌های بزرگ سخت است
\end{itemize}

\subsection{نیازها}
\begin{itemize}
    \item پنل ادمین ساده و فارسی
    \item افزودن محصول در کمتر از ۲ دقیقه
    \item مشاهدهٔ سفارش‌ها در یک صفحه
    \item تغییر وضعیت سفارش
\end{itemize}

\subsection{سناریوی استفاده}
\begin{quote}
صبح مغازه را باز می‌کند. با لپ‌تاپ وارد پنل ادمین می‌شود.
سه محصول جدید اضافه می‌کند. سفارش‌های دیشب را می‌بیند
و وضعیت دو سفارش را به «تحویل‌شده» تغییر می‌دهد.
\end{quote}

% ============================================================
\section{پرسونا ۲ — مریم (مشتری نهایی)}

\subsection{مشخصات دموگرافیک}
\begin{center}
\begin{tabular}{|l|l|}
\hline
\textbf{ویژگی} & \textbf{مقدار} \\
\hline
نام & مریم احمدی \\
\hline
سن & ۳۲ سال \\
\hline
جنسیت & زن \\
\hline
شغل & کارمند بانک \\
\hline
محل زندگی & تهران \\
\hline
تحصیلات & کارشناسی \\
\hline
مهارت فنی & بالا \\
\hline
وضعیت تأهل & متأهل، یک فرزند \\
\hline
\end{tabular}
\end{center}

\subsection{نقل‌قول}
\begin{quote}
«وقت ندارم به مغازه بروم. می‌خواهم سریع یخچال مناسب را پیدا کنم.»
\end{quote}

\subsection{اهداف}
\begin{itemize}
    \item خرید سریع و آسان لوازم خانگی
    \item مقایسهٔ محصولات
    \item اطمینان از اصالت و گارانتی
    \item تحویل در محل
\end{itemize}

\subsection{دردها}
\begin{itemize}
    \item وقت کافی ندارد
    \item سایت‌ها شلوغ و پرتبلیغ‌اند
    \item نمی‌داند کدام فروشنده قابل‌اعتماد است
    \item مقایسه سخت است
\end{itemize}

\subsection{نیازها}
\begin{itemize}
    \item جست‌وجو و فیلتر پیشرفته
    \item مقایسهٔ ۲-۳ محصول
    \item صفحهٔ محصول با گارانتی
    \item سبد خرید ساده
    \item ثبت سفارش سریع
\end{itemize}

\subsection{سناریوی استفاده}
\begin{quote}
در استراحت ناهار با موبایل وارد سایت می‌شود.
جست‌وجو می‌کند «یخچال ساید». با فیلتر قیمت
و برند، سه گزینه می‌بیند. دو تا را مقایسه می‌کند.
یکی را سفارش می‌دهد.
\end{quote}

% ============================================================
\section{پرسونا ۳ — رضا (انباردار)}

\subsection{مشخصات دموگرافیک}
\begin{center}
\begin{tabular}{|l|l|}
\hline
\textbf{ویژگی} & \textbf{مقدار} \\
\hline
نام & رضا کریمی \\
\hline
سن & ۲۸ سال \\
\hline
جنسیت & مرد \\
\hline
شغل & انباردار \\
\hline
محل زندگی & کاشان \\
\hline
تحصیلات & دیپلم \\
\hline
مهارت فنی & پایین \\
\hline
\end{tabular}
\end{center}

\subsection{نقل‌قول}
\begin{quote}
«با کاغذ و قلم کار کردن خطا زیاد دارد.»
\end{quote}

\subsection{اهداف}
\begin{itemize}
    \item ثبت دقیق موجودی
    \item اطلاع سریع از کمبود
    \item هماهنگی با فروش آنلاین
    \item کاهش خطای انسانی
\end{itemize}

\subsection{دردها}
\begin{itemize}
    \item سیستم کاغذی پرخطاست
    \item نمی‌داند چه در انبار است
    \item هماهنگی سخت است
    \item با نرم‌افزار راحت نیست
\end{itemize}

\subsection{نیازها}
\begin{itemize}
    \item رابط کاربری بسیار ساده
    \item ثبت موجودی با چند کلیک
    \item هشدار کمبود
    \item گزارش روزانه
\end{itemize}

\subsection{سناریوی استفاده}
\begin{quote}
صبح با تبلت وارد پنل انبار می‌شود.
۱۰ ماشین لباسشویی جدید را ثبت می‌کند.
لیست کمبودها را می‌بیند و به خانم رضایی خبر می‌دهد.
\end{quote}

% ============================================================
\section{ماتریس اولویت‌بندی}

\begin{center}
\begin{tabular}{|l|c|c|}
\hline
\textbf{پرسونا} & \textbf{تأثیر} & \textbf{اولویت} \\
\hline
مریم (مشتری) & بالا & اول \\
\hline
خانم رضایی (ادمین) & بالا & اول \\
\hline
رضا (انباردار) & متوسط & دوم \\
\hline
\end{tabular}
\end{center}

\subsection{نتیجه‌گیری}
\begin{itemize}
    \item \textbf{UI/UX} باید برای \textbf{مریم} (مشتری) بهینه باشد.
    \item \textbf{پنل ادمین} باید برای \textbf{خانم رضایی} ساده باشد.
    \item \textbf{پنل انبار} می‌تواند در Sprint 3 اضافه شود.
\end{itemize}

% ============================================================
\section{اتصال به User Stories}

پرسوناها پایهٔ User Storyها هستند:

\begin{center}
\begin{tabular}{|l|l|}
\hline
\textbf{پرسونا} & \textbf{نمونه User Story} \\
\hline
مریم & می‌خواهم محصولات را بر اساس قیمت فیلتر کنم \\
\hline
خانم رضایی & می‌خواهم محصول جدید اضافه کنم \\
\hline
رضا & می‌خواهم موجودی را به‌روز کنم \\
\hline
\end{tabular}
\end{center}

% ============================================================
\appendix
\section{تاریخچهٔ تغییرات}

\begin{center}
\begin{tabular}{|c|c|l|l|}
\hline
\textbf{نسخه} & \textbf{تاریخ} & \textbf{تغییرات} & \textbf{تهیه‌کننده} \\
\hline
۱.۰ & \today & نسخهٔ اولیه & ابراهیم فضلی \\
\hline
\end{tabular}
\end{center}

\end{document}
